The dominant electroweak effect on the $m_W$ measurement is final-state QED radiation~\cite{EWK},
which reduces the momentum of the charged lepton from the $W$-boson decay.
The experiments model this radiation using generators that resum multiple soft photon
emissions above an energy threshold. Uncertainties on the modelling of electroweak
corrections include:
(1) the perturbative calculation of photon radiation, including the modelling of single-photon and multi-photon
emission and the matching of the fixed-order and all-orders descriptions
(2) the energy threshold for producing final-state photons; and 
(3) higher-order corrections from final-state $e^+ e^-$ pair production.
Tables~\ref{tbl:ewkelectrons} and  \ref{tbl:ewkmuons} list the size of these uncertainties
for each experiment in the electron and muon channels, respectively.
The uncertainties are completely correlated between the decay channels.

\begin{table}[!tp]
  \centering 
  \begin{tabular}{lccc}
  \hline \hline 
     Uncertainty        & CDF       & D0     & ATLAS     \\ 
     \hline 
    Perturbative photon rad.   & 2.3 (2.3) & 5 (5) & 2.5 (3.3) \\ 
    Photon energy cutoff  & 1 (1)  & 2 (1) & ---    \\
    FSR $e^+e^-$       & 1 (1)     & ---  & 0.8 (3.6) \\
    \hline 
    Total              & 2.7 (2.7)     & 7 (7) & 2.6 (4.9) \\
    \hline \hline 
  \end{tabular}
  \caption{QED uncertainties in MeV on the $m_W$ measurement in the electron 
  channel using the \mT~(\pT) fit.  The uncertainties are uncorrelated except
  for those due to the perturbative photon radation calculation, which is taken
  to be 100\% correlated between D0 and ATLAS, and to the
  photon energy cutoff, taken to be 100\% correlated between CDF and D0. 
  }  
\label{tbl:ewkelectrons}
\end{table}

% LHCb: QED envelope pythia/photos/herwig (7 MeV) around average + powheg ew w vs w/o (5 MeV)
\begin{table}[!tp]
  \centering 
  \begin{tabular}{lccc}
  \hline \hline 
    Uncertainty          & CDF       & ATLAS     & LHCb \\
    \hline 
    Perturbative photon rad.   & 2.3 (2.3) & 2.5 (3.5) & 8.6 \\
    Photon energy cutoff & 1 (1)  & ---          & --- \\ 
    FSR $e^+e^-$         & 1 (1)     & 0.8 (3.6) & --- \\ 
    Total                & 2.7 (2.7) & 2.6 (5.6) & 8.6 \\
\hline \hline
  \end{tabular}
  \caption{QED uncertainties in MeV on the $m_W$ measurement in the muon 
  channel for ATLAS and CDF using the \mT~(\pT) fit, and for LHCb. 
  The uncertainties are taken to be uncorrelated between the experiments.
  }  
\label{tbl:ewkmuons}
\end{table}

To estimate the uncertainty from the limitations of the shower model relative to the 
matrix-element calculation, D0 and ATLAS perform a direct comparison between PHOTOS 
and WGRAD~\cite{Baur:1998kt,Baur:2004ig} or WINHAC~\cite{Placzek:2003zg,Placzek:2009jy,Placzek:2013moa}, 
respectively.  Since ATLAS and D0 use the same shower model, their uncertainties are 
considered as correlated.  LHCb estimates the uncertainty with a hybrid approach of 
comparing \textsc{Powheg} with and without the NLO EW calculation, and the range of 
the PHOTOS, \textsc{Pythia8}, and \textsc{Herwig} shower models.  The average of the 
measurements from the different shower models is taken as the central value, so the 
uncertainty is considered as uncorrelated.  CDF uses a third strategy, applying a 
correction to the measurement using the 
HORACE~\cite{CarloniCalame:2003ux,CarloniCalame:2006zq,CarloniCalame:2007cd} 
generator, which matches multiple-photon radiation to the $O(\alpha)$ calculation.
The residual uncertainties are largely due to MC statistics, and are considered as 
uncorrelated.

The shower model includes a lower threshold on the emitted photon energy, expressed 
as a ratio with respect to the energy of the lepton from the $W$ boson decay.  
CDF uses a threshold of $10^{-5}$ and determines the uncertainty by increasing 
the threshold by a factor of 3.  D0 uses a similar procedure except with an increase 
from $2.5 \times 10^{-4}$ to $2 \times 10^{-2}$.  These uncertainties are taken to be 
completely correlated.

To account for the higher-order process of an off-shell final-state photon 
splitting into an $e^+ e^-$ pair, CDF applies an effective radiator 
approximation to the radiated photons.  ATLAS does not apply a correction, 
instead taking the uncertainty from a PHOTOS model of this process.  
The uncertainties are treated as uncorrelated.
